\section{Economy}

\subsection{Salaries}

\subsection{Cost of Living}

\subsection{Stability}

\subsection{Pension Policy}

\subsection{Incentives for Foreigners}

\subsection{Tax}

\subsection{Public Transportation}

\subsection{Settling Permanently}

\subsection{Rating Guide}
When rating a country's economy, several key factors should be
considered to provide a comprehensive assessment:

\begin{enumerate}
  \item \textbf{Human Development Index (HDI)}: The Human Development
  Index is a composite index that takes into account three key
  dimensions of human well-being: life expectancy (health),
  education (knowledge), and per capita income (standard of living).
  HDI provides a more comprehensive picture of a country's overall
  development and welfare compared to GDP alone.
  
  \item \textbf{Housing Affordability Index}: the Housing Affordability
  Index assesses the ease with which individuals or families can afford
  to purchase a home. It considers factors such as median housing prices,
  household income, and interest rates to determine the affordability
  level of housing in a given area or country. A higher Housing
  Affordability Index score indicates that it is easier for people
  to buy their own homes.
  
  \item \textbf{Unemployment Rate}: The percentage of the labor force
  that is unemployed. Lower unemployment rates signify a healthier job
  market and better economic conditions.
  
  \item \textbf{Inflation Rate}: The rate at which the general price
  level of goods and services increases. Low and stable inflation is
  preferable for economic stability.
  
  \item \textbf{Trade Balance}: The difference between a country's
  exports and imports. A trade surplus (more exports) is generally
  desirable for economic strength.
  
  \item \textbf{Government Debt}: The total amount of money a country
  owes. Lower government debt is preferred for economic sustainability.
  
  \item \textbf{Fiscal Policy}: The government's revenue and spending
  decisions. Responsible fiscal policies support economic stability
  and growth.
  
  \item \textbf{Monetary Policy}: The central bank's management of the
  money supply and interest rates. Effective monetary policies
  influence economic growth and stability.
  
  \item \textbf{Infrastructure}: The quality of transportation,
  communication, and energy systems. Adequate infrastructure
  facilitates economic development.
  
  \item \textbf{Income Distribution}: How wealth and income are distributed
  among the population. A more equal income distribution can lead to
  social stability.
  
  \item \textbf{Technology and Innovation}: The level of technological
  advancement and investment in research and development. Innovation
  drives economic progress.
\end{enumerate}
